\begin{abstract}
Imagined-speech EEG classifiers report small margins above chance, so the evaluation protocol decides
what those margins mean. We audit the 39-way semantic-category benchmark released with Chisco, a
large public imagined-speech corpus, for participant S1, using the released code. Its stimulus
order is clustered by category, and its split places each evaluation trial between two training
trials, almost always from the same recording run. A classifier that never sees EEG and copies the label of the preceding
training trial scores 18.8\% on the released evaluation set, above the published 13.88\%, and
per-trial EEG summary statistics identify a held-out trial's recording run with 98.8\% accuracy.
Running the released script with its documented configuration on our hardware prints 9.69\%. Its
reporting rule, the best of several hundred evaluations on the evaluation set, gives 6.4\% to a model
trained on permuted labels. When whole runs are held out, real labels beat permuted labels by 0.64
points, within the variation between seeds, and a reading-task control shows that the pipeline has
limited sensitivity once recording-time cues are gone. The published number therefore does not, on its
own, establish semantic decoding. We recommend reporting position-only baselines on the exact split,
holding out whole runs, and selecting checkpoints on data that is never scored. Code, tests and
result files are available at \url{https://github.com/NoelSason/eegwords}.
\end{abstract}
